% GENERATED FILE. Edit canonical lessons, then run npm run generate:books.
% canonical-source-hash: fnv1a64:eecfeea175cad350
% canonical-lessons: FR-C450-caissier, FR-C450-encaisser, FR-C450-debiter, FR-C450-terminal, FR-C450-liquide

\chapter{À la caisse}
\label{ch:fr-a-la-caisse}
\hlchaptermodality{450}

\begin{chapteropening}
\textbf{By the end of this chapter:} \emph{I can pay at the till: the cashier, taking money, charging a card, the card machine and cash.}

Complete the last lesson of chapter 450: I can pay at the till: the cashier, taking money, charging a card, the card machine and cash.
\end{chapteropening}

\section[le caissier]{le caissier — the cashier}

\label{lesson:FR-C450-caissier}



\begin{quote}
Before the new one: say the French for to send off, then the French for to tie up.
\end{quote}

\subsection*{You'll want to know: le caissier}
\textbf{le caissier} — \textquotedblleft{}the cashier\textquotedblright{}.

\begin{quote}
Le caissier m'a rendu la monnaie.
\end{quote}

\begin{quote}
\emph{The cashier gave me my change.}
\end{quote}

Feminine \emph{la caissière}.

\begin{cousinweb}
\emph{caissier} is built on \emph{caisse} 'till, box' (EtymDB), which goes back to Latin \emph{capsa} 'box', probably through Provençal.
\end{cousinweb}

\begin{grammarlens}[title={What you've built}]
One more word for telling, explaining and describing at B1.
\end{grammarlens}

\subsection*{Your turn}
\begin{itemize}
\raggedright
  \item \emph{Say it:} \emph{le caissier}
  \item \emph{Say it:} \emph{le caissier}, once more
  \item \emph{Say it:} a sentence of your own with \emph{le caissier}
  \item \emph{Recall:} say the French for to send off, then the French for to tie up, then say \emph{le caissier} again
\end{itemize}

\begin{tcolorbox}[breakable,colback=teal!4,colframe=teal!35!black,title={Before you move on}]
What does le caissier mean? (\textquotedblleft{}The cashier\textquotedblright{}.) Say it once more.
\end{tcolorbox}
\section[encaisser]{encaisser — to cash in, to take (money); to take (a blow)}

\label{lesson:FR-C450-encaisser}



\begin{quote}
Before the new one: say the French for to tie up, then the French for the cashier.
\end{quote}

\subsection*{You'll want to know: encaisser}
\textbf{encaisser} — \textquotedblleft{}to cash in, to take (money); to take (a blow)\textquotedblright{}.

\begin{quote}
Le commerçant encaisse le paiement.
\end{quote}

\begin{quote}
\emph{The shopkeeper takes the payment.}
\end{quote}

Informal: \emph{je ne peux pas l'encaisser} — I cannot stand him.

\begin{cousinweb}
\emph{encaisser} is built from \emph{en-} and \emph{caisse} 'box, till', which goes back to Latin \emph{capsa} (Lewis \& Short): to put the money in the till.
\end{cousinweb}

\begin{grammarlens}[title={What you've built}]
One more word for telling, explaining and describing at B1.
\end{grammarlens}

\subsection*{Your turn}
\begin{itemize}
\raggedright
  \item \emph{Say it:} \emph{encaisser}
  \item \emph{Say it:} \emph{encaisser}, once more
  \item \emph{Say it:} a sentence of your own with \emph{encaisser}
  \item \emph{Recall:} say the French for to tie up, then the French for the cashier, then say \emph{encaisser} again
\end{itemize}

\begin{tcolorbox}[breakable,colback=teal!4,colframe=teal!35!black,title={Before you move on}]
What does encaisser mean? (\textquotedblleft{}To cash in, to take (money); to take (a blow)\textquotedblright{}.) Say it once more.
\end{tcolorbox}
\section[débiter]{débiter — to charge (an account); to cut up}

\label{lesson:FR-C450-debiter}



\begin{quote}
Before the new one: say the French for the cashier, then the French for to cash in.
\end{quote}

\subsection*{You'll want to know: débiter}
\textbf{débiter} — \textquotedblleft{}to charge (an account); to cut up\textquotedblright{}.

\begin{quote}
Le magasin a débité ma carte deux fois.
\end{quote}

\begin{quote}
\emph{The shop charged my card twice.}
\end{quote}

\emph{Le débit} — the debit; also the flow of a river.

\begin{cousinweb}
\emph{débiter} is built on \emph{débit} (EtymDB), from Latin \emph{dēbitum} 'what is owed', from \emph{dēbēre} 'to owe'.
\end{cousinweb}

\begin{grammarlens}[title={What you've built}]
One more word for telling, explaining and describing at B1.
\end{grammarlens}

\subsection*{Your turn}
\begin{itemize}
\raggedright
  \item \emph{Say it:} \emph{débiter}
  \item \emph{Say it:} \emph{débiter}, once more
  \item \emph{Say it:} a sentence of your own with \emph{débiter}
  \item \emph{Recall:} say the French for the cashier, then the French for to cash in, then say \emph{débiter} again
\end{itemize}

\begin{tcolorbox}[breakable,colback=teal!4,colframe=teal!35!black,title={Before you move on}]
What does débiter mean? (\textquotedblleft{}To charge (an account); to cut up\textquotedblright{}.) Say it once more.
\end{tcolorbox}
\section[le terminal]{le terminal — the card machine; the terminal}

\label{lesson:FR-C450-terminal}



\begin{quote}
Before the new one: say the French for to cash in, then the French for to charge.
\end{quote}

\subsection*{You'll want to know: le terminal}
\textbf{le terminal} — \textquotedblleft{}the card machine; the terminal\textquotedblright{}.

\begin{quote}
Elle pose sa carte sur le terminal.
\end{quote}

\begin{quote}
\emph{She taps her card on the card machine.}
\end{quote}

\emph{Le terminal de paiement}; also an airport terminal.

\begin{cousinweb}
\emph{terminal} is a learned borrowing from Late Latin \emph{terminālis} 'of a boundary' (EtymDB), from \emph{terminus} 'end, limit'.
\end{cousinweb}

\begin{grammarlens}[title={What you've built}]
One more word for telling, explaining and describing at B1.
\end{grammarlens}

\subsection*{Your turn}
\begin{itemize}
\raggedright
  \item \emph{Say it:} \emph{le terminal}
  \item \emph{Say it:} \emph{le terminal}, once more
  \item \emph{Say it:} a sentence of your own with \emph{le terminal}
  \item \emph{Recall:} say the French for to cash in, then the French for to charge, then say \emph{le terminal} again
\end{itemize}

\begin{tcolorbox}[breakable,colback=teal!4,colframe=teal!35!black,title={Before you move on}]
What does le terminal mean? (\textquotedblleft{}The card machine; the terminal\textquotedblright{}.) Say it once more.
\end{tcolorbox}
\section[le liquide]{le liquide — cash; the liquid}

\label{lesson:FR-C450-liquide}



\begin{quote}
Before the new one: say the French for to charge, then the French for the card machine.
\end{quote}

\subsection*{You'll want to know: le liquide}
\textbf{le liquide} — \textquotedblleft{}cash; the liquid\textquotedblright{}.

\begin{quote}
Le petit café accepte seulement le liquide.
\end{quote}

\begin{quote}
\emph{The little café only takes cash.}
\end{quote}

\emph{Payer en liquide} — to pay cash.

\begin{cousinweb}
\emph{liquide} is a learned borrowing from Latin \emph{liquidus} 'fluid' (EtymDB), from \emph{liquēre} 'to be fluid': money that flows freely.
\end{cousinweb}

\begin{grammarlens}[title={What you've built}]
One more word for telling, explaining and describing at B1.
\end{grammarlens}

\subsection*{Your turn}
\begin{itemize}
\raggedright
  \item \emph{Say it:} \emph{le liquide}
  \item \emph{Say it:} \emph{le liquide}, once more
  \item \emph{Say it:} a sentence of your own with \emph{le liquide}
  \item \emph{Recall:} say the French for to charge, then the French for the card machine, then say \emph{le liquide} again
\end{itemize}

\begin{tcolorbox}[breakable,colback=teal!4,colframe=teal!35!black,title={Before you move on}]
What does le liquide mean? (\textquotedblleft{}Cash; the liquid\textquotedblright{}.) Say it once more.
\end{tcolorbox}
